\vspace{2.5pt plus 3pt minus 0.5pt}\noindent\hspace*{1.2em}\kbvnum{(3)}\hspace{0.55em}ನಾನು ಪ್ರತಿ ದಿನ ರಾತ್ರಿ ಹತ್ತು ಗಂಟೆಗೆ ಮಲಗುತ್ತೇನೆ. ಈ ವಾಕ್ಯಗಳಲ್ಲಿ ದಪ್ಪಕ್ಷರಗಳಲ್ಲಿರುವ ಪದಗಳು ಸಂಖ್ಯಾವಾಚಕಗಳು.\par

\vspace{4pt plus 4pt minus 1pt}\noindent{\kbvsecfont\fontsize{11pt}{14.5pt}\selectfont\color{kbvaccentdark} (೪)\hspace{0.45em}ಸಂಖ್ಯೇಯ ವಾಚಕಗಳು :}\hspace{0.5em}"ಸಂಖ್ಯಾತ್ಮಕ ನಾಮಪದಗಳಿಗೆ ಸಂಖ್ಯೇಯ ವಾಚಕಗಳೆಂದು ಹೆಸರು. ಉದಾಹರಣೆಗೆ, “ದ್ರೌಪದಿಯು ಐವರು ಮಡದಿಯಾಗಿದ್ದಳು” ; ಈ ವಾಕ್ಯ ದಲ್ಲಿ ಐದು ಜನರು ಎಂಬ ಅರ್ಥದಲ್ಲಿ “ಐವರು” ಎಂಬ ಪದ ಬಳಸಲಾಗಿದೆ. ಇದು ಸಂಖ್ಯೇಯವಾಚಕ.\par

\vspace{2.5pt plus 3pt minus 0.5pt}\noindent\hspace*{1.2em}\kbvnum{(1)}\hspace{0.55em}ಓಟದ ಪಂದ್ಯದಲ್ಲಿ ಗೆದ್ದವರಲ್ಲಿ ನಾನು ಎರಡನೆಯವನು.\par

\vspace{2.5pt plus 3pt minus 0.5pt}\noindent\hspace*{1.2em}\kbvnum{(2)}\hspace{0.55em}ಕನ್ನಡದ ಶ್ರೇಷ್ಠ ಸಾಹಿತಿಗಳಲ್ಲಿ ಕುವೆಂಪು ಒಬ್ಬರು. ಈ ವಾಕ್ಯಗಳಲ್ಲಿ ದಪ್ಪಕ್ಷರಗಳಲ್ಲಿರುವ ಪದಗಳು ಸಂಖ್ಯೇಯವಾಚಕಗಳು.\par

\vspace{4pt plus 4pt minus 1pt}\noindent{\kbvsecfont\fontsize{11pt}{14.5pt}\selectfont\color{kbvaccentdark} (೫)\hspace{0.45em}ಪರಿಮಾಣವಾಚಕಗಳು :}\hspace{0.5em}ಅಂದಾಜಿನ ಮೇರೆಗೆ ವಸ್ತುಗಳ ಪರಿಮಾಣ ಸೂಚಿಸುವ ನಾಮಪದ ಗಳಿಗೆ ಪರಿಮಾಣ ವಾಚಕಗಳೆಂದು ಹೆಸರು.\par

\noindent\hspace*{1.2em}\kbvlabel{ಉದಾ}\hspace{0.6em}ಇಷ್ಟು, ಅಷ್ಟು, ರವೆಯಷ್ಟು, ಬೆಟ್ಟದಷ್ಟು, ಕೆಲವು, ಹಲವು—ಇವು ಪರಿಮಾಣವಾಚಕಗಳು.\par

\vspace{4pt plus 4pt minus 1pt}\noindent{\kbvsecfont\fontsize{11pt}{14.5pt}\selectfont\color{kbvaccentdark} (೬)\hspace{0.45em}ಪ್ರಕಾರವಾಚಕಗಳು :}\hspace{0.5em}ಹೋಲಿಕೆಯ ಅರ್ಥ `ಬರುವಂತೆ' ಮನುಷ್ಯ, ಪ್ರಾಣಿ, ಸ್ಥಳ, ವಸ್ತು ಇವುಗಳ ರೀತಿ ಸ್ವಭಾವ, ಸ್ಥಿತಿ ಯಾವ ಪ್ರಕಾರವಾಗಿರುವುದೆಂಬುದನ್ನು ತಿಳಿಸುವ ನಾಮಪದಗಳಿಗೆ ಪ್ರಕಾರವಾಚಕಗಳೆಂದು ಹೆಸರು. 'ಅಂಥ, ಇಂಥ, ಅಂತಹ, ಇಂತಹ ಮೊದಲಾದ ಪದಗಳನ್ನು ಪ್ರಕಾರವಾಚಕಗಳಾಗಿ ಬಳಸುವುದುಂಟು.\par

\noindent\hspace*{1.2em}\kbvlabel{ಉದಾ}\hspace{0.6em}ಅಂಥ ವ್ಯಕ್ತಿ ಮತ್ತೊಬ್ಬನಿಲ್ಲ.\par

\noindent\hspace*{2.2em}ಇಂಥ ಜನಕ್ಕೆ ಈ ಪಾಡು ಬರಬಾರದಿತ್ತು.\par

\noindent\hspace*{2.2em}ಅಂತಹವನು ಹುಟ್ಟುವುದು ಅಪರೂಪ.\par

\noindent\hspace*{2.2em}ನೀನು ಇಂತಹ ಕೆಲಸ ಮಾಡಿದ್ದು ತಪ್ಪು.\par
